\section{A uniform residue-two Schur bound}\label{app:r2}

We prove Theorem~\ref{thm:r2}. In this appendix $n=8k+2=4\ell+2$,
where $\ell=2k$ is even and $k\ge6$. The letter $\ell$ counts four-vertex
blocks and is distinct from the cell count $m$ in the main text. Put
\[
 M_n=\frac{198}{25}I_n-B_n^2.
\]
All matrix inequalities in this appendix concern real symmetric matrices;
$X\succ0$ means positive definite. The norm $\norm{\cdot}_2$ is the
Euclidean operator norm, and $\norm{\cdot}_F$ is the Frobenius norm.

\subsection{The exact block recurrence}

Partition the vertices into
\[
 V_0=\{0,1\},\qquad
 V_j=\{2+4(j-1),\ldots,5+4(j-1)\}\quad(1\le j\le\ell).
\]
The diagonal block on each $V_j$, $j\ge1$, is
\[
 D=\begin{pmatrix}
 98/25&0&-1&0\\0&98/25&0&-1\\-1&0&98/25&0\\0&-1&0&98/25
 \end{pmatrix}.
\]
The block $M_n[V_j,V_{j+1}]$ equals $E_+$ for odd $j$ and $E_-$ for
even $j$, where
\[
 E_+=\begin{pmatrix}-1&0&0&0\\0&1&0&0\\-1&2&1&0\\2&-1&0&-1\end{pmatrix},
 \qquad
 E_-=\begin{pmatrix}-1&0&0&0\\0&1&0&0\\-1&-2&1&0\\-2&-1&0&-1\end{pmatrix}.
\]
The exceptional blocks are
\begin{gather*}
 G_0=M_n[V_0,V_0]=(98/25)I_2,\qquad H_0=D,\\
 R_0=M_n[V_0,V_1]=\begin{pmatrix}-1&-2&1&0\\-2&-1&0&-1\end{pmatrix},\\
 C_0=M_n[V_0,V_\ell]=\begin{pmatrix}-1&0&-1&0\\0&-1&0&-1\end{pmatrix},
 \qquad
 W_0=M_n[V_1,V_\ell]=\begin{pmatrix}0&0&-1&0\\0&0&0&1\\0&0&0&0\\0&0&0&0\end{pmatrix}.
\end{gather*}
All other nonadjacent bulk blocks vanish. To verify these identities at
all lengths, the only nonzero entries of $B_n^2$ on the forward cyclic
diagonals, apart from transposes, are
\begin{align*}
 (B_n^2)_{i,i}&=4,&
 (B_n^2)_{i,i+1}&=\tau_{i-1}+\tau_i,&
 (B_n^2)_{i,i+2}&=1,\\
 (B_n^2)_{i,i+3}&=\tau_i+\tau_{i+1},&
 (B_n^2)_{i,i+4}&=\tau_i\tau_{i+2}.&&
\end{align*}
They follow by enumerating the possible two-edge walks. Since $n\ge50$,
these cyclic offsets do not collide. Substitution of $t^k\|(1,-1)$ gives
the displayed blocks without a length-dependent assumption.

Let $E_j=E_+$ for even $j$ and $E_j=E_-$ for odd $j$. Set $X_0=D$.
Before closing the right boundary, elimination of the current four-vertex
pivot gives the length-independent recurrence
\begin{align}
 X_{j+1}&=D-E_j^TX_j^{-1}E_j,&
 R_{j+1}&=-R_jX_j^{-1}E_j,\notag\\
 W_{j+1}&=-E_j^TX_j^{-1}W_j,&
 G_{j+1}&=G_j-R_jX_j^{-1}R_j^T,\label{eq:recurrence}\\
 H_{j+1}&=H_j-W_j^TX_j^{-1}W_j,&
 C_{j+1}&=C_j-R_jX_j^{-1}W_j.\notag
\end{align}
The terminal elimination index is $p=\ell-2$, which is even.
At that step only, the coupling to the retained last block becomes
$W_p+E_+$. Thus the normalized six-dimensional core is
\begin{equation}
 S_\ell=
 \begin{pmatrix}
 G_p-R_pX_p^{-1}R_p^T&C_p-R_pX_p^{-1}(W_p+E_+)\\
 *&H_p-(W_p+E_+)^TX_p^{-1}(W_p+E_+)
 \end{pmatrix}.
 \label{eq:core}
\end{equation}
Here $*$ denotes the transpose block. Repeated Schur congruence proves
that $M_n\succ0$ if all eliminated pivots and $S_\ell$ are positive.
The subscript of $S_\ell$ is a block count: in particular, $S_{26}$
corresponds to graph order $106$.

\subsection{Finite exact premises}

Define $F_\pm(X)=D-E_\pm^TX^{-1}E_\pm$ and
$\Phi=F_-\circ F_+$. Put $Z=\Phi^{12}(D)=X_{24}$,
$Y=F_+(Z)$, and $L_Z=Z^{-1}E_+Y^{-1}E_-$. Use the rational matrices
\begin{align*}
 P&=\frac1{10000}\begin{pmatrix}
10766&87&19&974\\87&12664&148&-2418\\19&148&10093&-25\\974&-2418&-25&14009
\end{pmatrix},\\
 Q&=\frac1{10000}\begin{pmatrix}
11503&614&990&-1101\\614&10470&15&113\\990&15&12299&-2632\\-1101&113&-2632&13260
\end{pmatrix}.
\end{align*}
Set $r=10^{-10}$. The following finite inequalities are the complete
certificate premises used below:
\begin{enumerate}
\item $X_0,\ldots,X_{24}\succ0$ and $Z,Y\succ I_4/2$;
\item $(9/10)I_4\prec P,Q\prec2I_4$;
\item $L_Z^TPL_Z\prec(2/5)P$ and $L_ZQL_Z^T\prec(2/5)Q$;
\item $\norm{\Phi(Z)-Z}_F<r/40$;
\item $R_{24}QR_{24}^T\prec10^{-10}I_2$ and
$W_{24}^TQW_{24}\prec10^{-10}I_4$;
\item $S_{26}\succ(1/50)I_6$;
\item $S_\ell\succ0$ for $\ell=12,14,16,18,20,22,24$.
\end{enumerate}
They have been verified using exact rational arithmetic. For clarity, their
verification requires no limiting or approximate matrix: all matrices are
specified by the displayed rational initial data and a finite number of
recurrence steps. Positive definiteness is checked by successive scalar
Schur pivots. For a symmetric matrix $K$, start with $K^{(0)}=K$; at each
step require $d_j=K^{(j)}_{11}>0$ and replace the remaining block by
$K^{(j+1)}=K^{(j)}_{22}-K^{(j)}_{21}(d_j)^{-1}K^{(j)}_{12}$.
All resulting scalars are rational, so each test is exact. The residual
test in item 4 is the rational inequality
\[
 \sum_{a,b}(\Phi(Z)-Z)_{ab}^2<(r/40)^2.
\]
The accompanying certificate records the matrices and the positive pivots.
An independent replay constructs $B_{106}$ from its edges, eliminates
$96$ scalar vertices to recover the entrance data, and verifies items
1--5 without generating that entrance by the block recurrence. Direct
scalar elimination from the graphs also verifies the seven base orders
$50,58,66,74,82,90,98$ and the normalized seed at $106$.

\subsection{A local Riccati contraction}

For symmetric $U$, define the order-unit norm
\[
 |U|_P=\inf\{a\ge0:-aP\preceq U\preceq aP\}
      =\norm{P^{-1/2}UP^{-1/2}}_2.
\]
Consider the closed ball
$\mathcal K=\{X=X^T:|X-Z|_P\le r\}$. Since $P\prec2I$,
\[
 \norm{X-Z}_2\le2r=:e.
\]
The certified lower bound $Z\succ I/2$ gives $X\succ I/3$ and
$\norm{X^{-1}}_2\le3$. The inverse identity yields
\[
 \norm{X^{-1}-Z^{-1}}_2\le6e.
\]
Both $E_\pm$ have squared Frobenius norm $14$. Writing $Y_X=F_+(X)$,
we obtain
\[
 \norm{Y_X-Y}_2\le84e<1/6,
 \qquad Y_X\succ I/3,
 \qquad \norm{Y_X^{-1}-Y^{-1}}_2\le504e.
\]
Define $L(X)=X^{-1}E_+Y_X^{-1}E_-$. Splitting its difference from
$L_Z$ into two terms gives
\begin{equation}
 \norm{L(X)-L_Z}_2
 \le14(6\cdot3+2\cdot504)e=14364e.
 \label{eq:Lperturb}
\end{equation}
For the two transfer norms
\[
 \norm{L}_{P,\mathrm{col}}=\norm{P^{1/2}LP^{-1/2}}_2,
 \qquad
 \norm{L}_{Q,\mathrm{row}}=\norm{Q^{-1/2}LQ^{1/2}}_2,
\]
conversion from the Euclidean norm costs at most
$\sqrt{20/9}<3/2$. Therefore the perturbation in either norm is less
than $21546e=43092r<10^{-4}$. The certified bounds at $Z$, together with
$\sqrt{2/5}+10^{-4}<2/3$, imply throughout $\mathcal K$ that
\begin{equation}
 L(X)^TPL(X)\prec\frac49P,
 \qquad L(X)QL(X)^T\prec\frac49Q.
 \label{eq:dual-contractions}
\end{equation}
Direct differentiation gives
$D\Phi_X[U]=L(X)^TUL(X)$. If $-aP\preceq U\preceq aP$, the first
inequality in \eqref{eq:dual-contractions} bounds this derivative in
$|\cdot|_P$ by $4a/9$. Integrating along a line segment in the convex
ball proves that $\Phi$ is $4/9$-Lipschitz there.

The residual premise gives
$|\Phi(Z)-Z|_P<(10/9)(r/40)=r/36$. Consequently
\[
 |\Phi(X)-Z|_P<r/36+(4/9)r=(17/36)r<r.
\]
The ball is invariant. Banach's fixed-point theorem supplies a unique
$X_*\in\mathcal K$ with $\Phi(X_*)=X_*$. The actual even trajectory
starts at $Z$, so, with $\theta=4/9$,
\begin{equation}
 \norm{X_{24+2h}-X_*}_2\le4r\theta^h
 \qquad(h\ge0).
 \label{eq:pivot-decay}
\end{equation}
The constant $4r$ follows already from the diameter of the ball. All these
even pivots and their intermediate odd partners exceed $I/3$. Together
with the finite entrance this proves positivity of every pivot needed
at any length.

\subsection{Response decay}

Two consecutive updates in \eqref{eq:recurrence} act by
$R\mapsto RL(X)$ and $W\mapsto L(X)^TW$. For these two orientations,
respectively, use
\[
 \alpha(R)=\norm{RQ^{1/2}}_2,
 \qquad \beta(W)=\norm{Q^{1/2}W}_2.
\]
The second inequality in \eqref{eq:dual-contractions} contracts both
quantities by $q=2/3$. Hence the entrance premises imply
\begin{align*}
 R_{24+2h}QR_{24+2h}^T&\prec10^{-10}q^{2h}I_2,\\
 W_{24+2h}^TQW_{24+2h}&\prec10^{-10}q^{2h}I_4.
\end{align*}
Since $Q\succ(9/10)I$, it follows that
\begin{equation}
 \norm{R_{24+2h}}_2,\ \norm{W_{24+2h}}_2<a q^h,
 \qquad a=1/30000.
 \label{eq:even-response}
\end{equation}
Each one-step transfer has norm at most $3\sqrt{14}<12$. Thus both
members of the $h$th pair obey
\begin{equation}
 \norm{R_j}_2,\ \norm{W_j}_2\le bq^h,
 \qquad b=12a=1/2500,
 \quad j\in\{24+2h,25+2h\}.
 \label{eq:pair-response}
\end{equation}
The row orientation $LQL^T$ is needed here. A bound on $L^TQL$ alone
would not justify these response estimates.

\subsection{The limiting core and its uniform tail}

Use the actual infinite trajectory, not a trajectory with constant pivots,
to define the absolutely convergent series
\begin{align*}
 G_\infty&=G_0-\sum_{j\ge0}R_jX_j^{-1}R_j^T,\\
 H_\infty&=H_0-\sum_{j\ge0}W_j^TX_j^{-1}W_j,\\
 C_\infty&=C_0-\sum_{j\ge0}R_jX_j^{-1}W_j.
\end{align*}
Convergence follows from \eqref{eq:pair-response} and
$\norm{X_j^{-1}}_2\le3$ after $j=24$. Set
\[
 S_\infty=\begin{pmatrix}
 G_\infty&C_\infty\\
 C_\infty^T&H_\infty-E_+^TX_*^{-1}E_+
 \end{pmatrix}.
\]
Take even $\ell\ge26$ and write its terminal index as
$p=\ell-2=24+2h$. Each of the three Schur-series suffixes after index
$p$ is bounded in norm by
\begin{equation}
 T_h=\frac{6b^2q^{2h}}{1-q^2}.
 \label{eq:Th}
\end{equation}
Indeed, every increment in the $j$th pair has norm at most
$3b^2q^{2j}$, and there are two increments per pair. The bound in
\eqref{eq:Th} harmlessly includes the term at $p$ as well.

The terminal linear terms from \eqref{eq:core} satisfy
\begin{align*}
 \norm{R_pX_p^{-1}E_+}_2&\le12a q^h,\\
 \norm{W_p^TX_p^{-1}E_++E_+^TX_p^{-1}W_p}_2&\le24a q^h.
\end{align*}
The inverse identity and \eqref{eq:pivot-decay}, with $\delta=4r$,
give
\[
 \norm{X_p^{-1}-X_*^{-1}}_2\le9\delta\theta^h,
 \qquad
 \norm{E_+^T(X_p^{-1}-X_*^{-1})E_+}_2
 \le144\delta\theta^h.
\]
For the differences of the three blocks in $S_\ell-S_\infty$ we
therefore have
\begin{align*}
 \norm{\Delta G}_2&\le T_h,\\
 \norm{\Delta C}_2&\le T_h+12a q^h,\\
 \norm{\Delta H}_2&\le T_h+24a q^h+144\delta\theta^h.
\end{align*}
The block estimate
\[
 \norm{\begin{pmatrix}U&V\\V^T&Z\end{pmatrix}}_2
 \le\norm U_2+2\norm V_2+\norm Z_2
\]
now gives the explicit uniform tail bound
\begin{align}
 \norm{S_\ell-S_\infty}_2
 &\le4T_h+48a q^h+144\delta\theta^h\notag\\
 &\le\frac{251089}{156250000}<\frac1{500}.
 \label{eq:tail}
\end{align}
The rational middle expression is the value of this majorant at $h=0$;
each geometric factor decreases for $h\ge0$. Because the series were
defined along the actual trajectory, the common initial terms cancel
exactly. Only the terminal pure-pivot term requires the inverse-limit
comparison above.

\subsection{Completion of the spectral estimate}

The exact seed $S_{26}\succ I_6/50$ and two applications of
\eqref{eq:tail} imply, for every even $\ell\ge26$,
\[
 \norm{S_\ell-S_{26}}_2<2/500=1/250,
 \qquad S_\ell\succ\left(\frac1{50}-\frac1{250}\right)I_6
                  =\frac2{125}I_6.
\]
Both tail errors are required: the seed is a finite core, not the limiting
core. The remaining even block counts $12\le\ell<26$ are exactly the
seven certified bases. Thus every core and every eliminated pivot is
positive. Schur congruence proves $M_n\succ0$, and the spectral theorem
then gives $\rho(B_n)^2<198/25$.

This use of Schur congruence asserts positivity only. The lower bound
$2/125$ for the reduced core is not a lower bound for the least Euclidean
eigenvalue of the original matrix $M_n$.
Finally, $\cos x>1-x^2/2$ for $x>0$ and $\pi^2<10$ yield
\[
 \tw(n)^2>8-\frac{20\pi^2}{n^2}
          >8-\frac{200}{n^2}\ge\frac{198}{25}
 \qquad(n\ge50).
\]
This completes the proof of Theorem~\ref{thm:r2}.

\paragraph{Reproducibility.}
\begingroup\raggedright
The accompanying source package includes
\nolinkurl{certificates/analytic/verify_r2_certificate.py} and
\nolinkurl{r2_exact_certificate.json}, together with the independent
\nolinkurl{replay_direct_graph_seed.py} and
\nolinkurl{replay_local_from_direct_graph.py} under
\nolinkurl{certificates/audit/}. For the separate antiperiodic check, the script
\nolinkurl{verify_antiperiodic_counterexample.py} under
\nolinkurl{certificates/new_conjecture/} constructs the graph,
checks the symbolic identities, and verifies all recorded integer minors.
The package README lists the exact execution commands and dependencies.
\par\endgroup
